% Kompilieren mit: TEXINPUTS=minted/source: xelatex -shell-escape %
\documentclass[12pt,compress,ngerman,utf8,t]{beamer}
\usepackage[ngerman]{babel}
\usepackage{comment}
\usepackage{minted}
\usepackage{wasysym}
\usepackage[all]{xy}
\usepackage[protrusion=true,expansion=false]{microtype}

\DeclareSymbolFont{extraup}{U}{zavm}{m}{n}
\DeclareMathSymbol{\varheart}{\mathalpha}{extraup}{86}
\DeclareMathSymbol{\vardiamond}{\mathalpha}{extraup}{87}

\title[Freie Monaden]{\smiley{} Freie Monaden \smiley}
\author[Augsburger Curry Club]{
  \includegraphics[scale=0.1]{images/a-monad-is-just} \\\ \\
  Ingo Blechschmidt \\[0.1em] \scriptsize\texttt{<iblech@speicherleck.de>}}
\date[2015-09-10]{\vspace*{-1.5em}\scriptsize Augsburger Curry-Club \\ 10. September 2015}

\usetheme{Warsaw}

\useinnertheme{rectangles}

\usecolortheme{seahorse}
\definecolor{mypurple}{RGB}{150,0,255}
\setbeamercolor{structure}{fg=mypurple}

\usefonttheme{serif}
\usepackage[T1]{fontenc}
\usepackage{libertine}

\renewcommand{\C}{\mathcal{C}}
\newcommand{\D}{\mathcal{D}}
\newcommand{\id}{\mathrm{id}}
\newcommand{\Id}{\mathrm{Id}}
\newcommand{\Hask}{\mathrm{Hask}}

\setbeamertemplate{navigation symbols}{}
\setbeamertemplate{headline}{}

\setbeamertemplate{title page}[default][colsep=-1bp,rounded=false,shadow=false]
\setbeamertemplate{frametitle}[default][colsep=-2bp,rounded=false,shadow=false,center]

\newcommand*\oldmacro{}%
\let\oldmacro\insertshorttitle%
\renewcommand*\insertshorttitle{%
  \oldmacro\hfill\insertframenumber\,/\,\inserttotalframenumber\hfill}

\newcommand{\hil}[1]{{\usebeamercolor[fg]{item}{\textbf{#1}}}}
\setbeamertemplate{frametitle}{%
  \vskip1em%
  \leavevmode%
  \begin{beamercolorbox}[dp=1ex,center]{}%
      \usebeamercolor[fg]{item}{\textbf{\textsf{\Large \insertframetitle}}}
  \end{beamercolorbox}%
}

\setbeamertemplate{footline}{%
  \leavevmode%
  \hfill%
  \begin{beamercolorbox}[ht=2.25ex,dp=1ex,right]{}%
    \usebeamerfont{date in head/foot}
    \insertframenumber\,/\,\inserttotalframenumber\hspace*{1ex}
  \end{beamercolorbox}%
  \vskip0pt%
}

\setbeameroption{show notes}
\setbeamertemplate{note page}[plain]

\begin{document}

\frame{\titlepage}

\section{Monoide}

\subsection{Definition und Beispiele}

\begin{frame}[fragile]\frametitle{Monoide}
  Ein \hil{Monoid} besteht aus
  \begin{itemize}
    \item einer Menge~$M$,
    \item einer Abbildung~$({\circ}) : M \times M \to M$ und
    \item einem ausgezeichneten Element~$1 \in M$,
  \end{itemize}
  sodass die \hil{Monoidaxiome} gelten: Für alle~$x,y,z \in M$
  \begin{itemize}
    \item $x \circ (y \circ z) = (x \circ y) \circ z$,
    \item $1 \circ x = x$,
    \item $x \circ 1 = x$.
  \end{itemize}

  \visible<2>{\begin{center}
    \hil{Beispiele:} \\
    natürliche Zahlen, Listen, Endomorphismen,
    Matrizen, \ldots
    \\[1em]
    \hil{Nichtbeispiele:} \\
    natürliche Zahlen mit Subtraktion,
    nichtleere Listen, \ldots
  \end{center}}

  \vspace*{-22em}
  \begin{columns}
    \begin{column}{0.65\textwidth}
    \end{column}
    \begin{column}{0.4\textwidth}
      \scriptsize\begin{block}{}
      \begin{minted}{haskell}
class Monoid m where
    (<>) :: m -> m -> m
    unit :: m
      \end{minted}
      \end{block}
    \end{column}
  \end{columns}
\end{frame}

\begin{frame}\frametitle{Axiome in Diagrammform}
  \[ \xymatrixcolsep{4pc}\xymatrixrowsep{4pc}\xymatrix{
    M \times M \times M \ar[r]^{\id \times ({\circ})} \ar[d]_{({\circ}) \times \id} & M \times M
    \ar[d]^{({\circ})} \\
    M \times M \ar[r]_{({\circ})} & M
  } \]
  \medskip
  \[ \xymatrixcolsep{4pc}\xymatrixrowsep{4pc}\xymatrix{
    & M \\
    M \ar[r]\ar[ru] & M \times M \ar[u] & \ar[l]\ar[lu] M
  } \]
\end{frame}

\begin{frame}[fragile]\frametitle{Monoidhomomorphismen}
  Eine Abbildung~$\varphi : M \to N$ zwischen Monoiden heißt genau dann
  \hil{Monoidhomomorphismus}, wenn
  \begin{itemize}
    \item $\varphi(1) = 1$ und
    \item $\varphi(x \circ y) = \varphi(x) \circ \varphi(y)$ für alle~$x,y \in M$.
  \end{itemize}

  \begin{center}
    \hil{Beispiele:} \\
    \mintinline{haskell}{length :: [a] -> Int} \\
    \mintinline{haskell}{sum :: [Int] -> Int} \\[1em]
    \hil{Nichtbeispiele:} \\
    \mintinline{haskell}{reverse :: [a] -> [a]} \\
    \mintinline{haskell}{head :: [a] -> a} \\
  \end{center}
\end{frame}

\subsection{Nutzen}

\begin{frame}\frametitle{Wozu?}
  \begin{itemize}
    \item Allgegenwärtigkeit
    \item Gemeinsamkeite und Unterschiede
    \item Generische Beweise
    \item Generische Algorithmen
  \end{itemize}
\end{frame}

\note{
  \begin{itemize}
    \item Monoide gibt es überall.
    \item Das Monoidkonzept hilft, Gemeinsamkeiten und Unterschiede
    zu erkennen und wertschätzen zu können.
    \item Man kann generische Beweise für beliebige Monoide führen.
    \item Man kann generische Algorithmen mit beliebigen Monoiden basteln.
  \end{itemize}
}


\subsection{Freie Monoide}

\begin{frame}\frametitle{Freie Monoide}
  Gegeben eine Menge~$X$ ohne weitere Struktur.
  Wie können wir auf möglichst ökonomische Art und Weise daraus einen
  Monoid~$F(X)$ gewinnen?
  \medskip
  \pause

  Spoiler: Der \hil{freie Monoid}~$F(X)$ auf~$X$ ist der Monoid der
  \hil{Listen} mit Elementen aus~$X$.
  \medskip
  \pause
  % auf der Konsole: Erzeuger passen und Relationen passen.

  Essenz der Freiheit: Jede beliebige Abbildung~$X \to M$ in einen Monoid~$M$
  stiftet genau einen Homomorphismus $F(X) \to M$.
  \[ \xymatrix{
    X \ar[rr] \ar[rd] && M \\
    & F(X) \ar@{-->}[ru]
  } \]
  \pause

  Der Funktor~$F : \mathrm{Set} \to \mathrm{Mon}$ ist \hil{linksadjungiert} zum
  Vergissfunktor~$\mathrm{Mon} \to \mathrm{Set}$.
\end{frame}

\note{
  In~$F(X)$ sollen neben den Elementen aus~$X$ nur so viele
  weitere Elemente enthalten sein, sodass man eine Verknüpfung~$({\circ})$ und
  ein Einselement definieren kann.

  In~$F(X)$ sollen nur die Rechenregeln gelten, die von den Monoidaxiomen
  gefordert werden.

  Die universelle Eigenschaft kann man schön in Haskell demonstrieren.
  \inputminted{haskell}{images/universal-property-free-monoid.hs}

  Wenn man tiefer in das Thema einsteigt, erkennt man, dass endliche Listen
  noch nicht der Weisheit letzter Schluss sind:
  \url{http://comonad.com/reader/2015/free-monoids-in-haskell/}
}


\section{Funktoren als Container}

\begin{frame}[fragile]\frametitle{Funktoren}
  Ein \hil{Funktor}~$F : \C \to \D$ zwischen Kategorien~$\C$ und~$\D$ ordnet
  \begin{itemize}
    \item jedem Objekt~$X \in \C$ ein Objekt~$F(X) \in \D$ und
    \item jedem Morphismus~$f : X \to Y$ in~$\C$ ein Morphismus~$F(f) : F(X)
    \to F(Y)$ in~$\D$ zu,
  \end{itemize}
  sodass die \hil{Funktoraxiome} erfüllt sind:
  \begin{itemize}
    \item $F(\id_X) = \id_{F(X)}$,
    \item $F(f \circ g) = F(f) \circ F(g)$.
  \end{itemize}
  \vfill

  In Haskell kommen Funktoren~$\Hask \to \Hask$ vor:
  \begin{minted}{haskell}
class Functor f where
    fmap :: (a -> b) -> (f a -> f b)
  \end{minted}
\end{frame}

\begin{frame}[fragile]\frametitle{Beispiele für Funktoren}
  \vspace*{-1em}
  \small
  \begin{minted}{haskell}
class Functor f where
    fmap :: (a -> b) -> (f a -> f b)

instance Functor [] where fmap f = map f

data Maybe a = Nothing | Just a
instance Functor Maybe where
    fmap f Nothing  = Nothing
    fmap f (Just x) = Just (f x)

data Id a = MkId a
instance Functor Id where
    fmap f (MkId x) = MkId (f x)

data Pair a = MkPair a a
instance Functor Pair where
    fmap f (MkPair x y) = MkPair (f x) (f y)
  \end{minted}
\end{frame}

\begin{frame}[fragile]\frametitle{Funktoren als Container}
  Ist~\texttt{f} ein Funktor, so stellen wir uns den
  Typ~\texttt{f a} als einen Typ von Containern von Werten vom
  Typ~\texttt{a} vor.

  Je nach Funktor haben die Container eine andere Form.
\end{frame}

\note{
  Die Vorstellung ist aus folgendem Grund plausibel: Aus einer
  Funktion~\texttt{a -> b} können wir mit \texttt{fmap} eine Funktion~\texttt{f
  a -> f b} machen. Also stecken wohl in einem Wert vom Typ~\texttt{f a}
  irgendwelche Werte vom Typ~\texttt{a}, die mit der gelifteten Funktion dann
  in Werte vom Typ~\texttt{b} umgewandelt werden.
}


\section{Monaden}

\begin{frame}[fragile]\frametitle{Monaden}
  Eine \hil{Monade} besteht aus
  \begin{itemize}
    \item einem Funktor~$M$,
    \item einer natürlichen Transformation~$M \circ M \Rightarrow M$ und
    \item einer natürlichen Transformation~$\Id \Rightarrow M$,
  \end{itemize}
  sodass die \hil{Monadenaxiome} gelten.
  \pause

  \begin{minted}{haskell}
class (Functor m) => Monad m where
    join   :: m (m a) -> m a
    return :: a -> m a
  \end{minted}

  \begin{columns}
    \begin{column}{0.42\textwidth}
      \begin{block}{Listen}
        \scriptsize
        \begin{minted}{haskell}
concat    :: [[a]] -> [a]
singleton :: a     -> [a]
        \end{minted}
      \end{block}
    \end{column}
    \begin{column}{0.62\textwidth}
      \begin{block}{Maybe}
        \scriptsize
        \begin{minted}{haskell}
join :: Maybe (Maybe a) -> Maybe a
Just :: a -> Maybe a
        \end{minted}
      \end{block}
    \end{column}
  \end{columns}
  \medskip
  Kein Beispiel: \mintinline{haskell}{Pair} von eben.
\end{frame}

\begin{frame}[fragile]\frametitle{Weitere Beispiele}
  \begin{minted}{haskell}
class (Functor m) => Monad m where
    join   :: m (m a) -> m a
    return :: a -> m a
  \end{minted}

  \begin{block}{Reader}
    \scriptsize
    \begin{minted}{haskell}
type Reader env a = env -> a

instance Functor Reader where
    fmap f k = f . k

instance Monad Reader where
    return x = \_   -> x
    join   k = \env -> k env env
    \end{minted}
  \end{block}

  \begin{block}{State}
    \scriptsize
    \begin{minted}{haskell}
type State s a = s -> (a,s)

instance Monad State where
    return x = \s -> (x,s)
    join   k = \s -> let (k',s') = k s in k' s'
    \end{minted}
  \end{block}
\end{frame}

\begin{frame}\frametitle{Die Monadenaxiome}
  {\scriptsize
  Sprechweise: Ein Wert vom Typ~\texttt{m (m (m a))} ist ein (äußerer) Container von
  (inneren) Containern von (ganz inneren) Containern von Werten vom
  Typ~\texttt{a}.\par}

  \[ \xymatrixcolsep{4pc}\xymatrixrowsep{4pc}\xymatrix{
    M \circ M \circ M \ar[r]^{\text{innen join}} \ar[d]_{\text{außen join}} & M \circ M
    \ar[d]^{\text{join}} \\
    M \times M \ar[r]_{\text{join}} & M
  } \]
  \medskip
  \[ \xymatrixcolsep{4pc}\xymatrixrowsep{4pc}\xymatrix{
    & M \\
    M \ar[r]_{\text{return}}\ar[ru] & M \circ M \ar[u]_{\text{join}} & \ar[l]^{\text{innen
    return}}\ar[lu] M
  } \]
\end{frame}

\end{document}

Monoide
* Definition und Beispiele
* Nutzen: Gemeinsamkeiten, generische Algorithmen, ...
* Freie Monoide: Ökonomieprinzip, universelle Eigenschaft

Funktoren als Container

Monaden
* Monaden als Container, die eine Art "join" unterstützen
* Monadenaxiome (mit Gegenüberstellung Set vs. End(Hask))

Freie Monaden
* Herleitung/Definition
* Universelle Eigenschaft
* Ausblick: Performanceprobleme
